\chapter{Glossary}
\label{app:glossary}

\begin{cctable}{Term}{Meaning}{1.35in}{3.47in}

\textbf{Address} & The name of a cell: its column letter followed by its row
number, as in \cellref{B9}. \\

\textbf{Banked RAM} & Memory reached 8K at a time through a window in the
address space, rather than directly. The whole workbook lives here. See
Appendix~\ref{app:memory}. \\

\textbf{Boolean} & A cell holding \msg{TRUE} or \msg{FALSE}. Counts as 1 or 0
in arithmetic. \\

\textbf{Cell} & One box of the worksheet. Holds a number, text, a date, a
boolean, a formula or an error value --- or nothing. \\

\textbf{Cell cursor} & The reverse-video block showing which cell the keyboard
is aimed at. \\

\textbf{Clipboard} & The one cell held by \menupath{Copy} for
\menupath{Paste}. \\

\textbf{Chart} & A picture of the numbers in the cursor's column --- bars, a
line or a pie --- drawn on the whole screen. Not stored in the workbook, but
\key{S} writes it to the card as \fname{CHART.BMP}.
Chapter~\ref{ch:charting}. \\

\textbf{CMDR-DOS} & The disk operating system of the Commander X16, which
\ccprod{} uses for all file access. \\

\textbf{Column} & A vertical line of cells, named \cellref{A} to
\cellref{IV}. \\

\textbf{Commit} & To finish an entry so it is stored, with \key{Return},
\key{Tab}, \keyc{Shift}{Tab} or an up or down arrow. \\

\textbf{CSV} & Comma-separated values: plain text, one line per row, fields
separated by commas. \\

\textbf{Current cell} & The cell the cursor is on. Every one-cell command acts
on it. \\

\textbf{DEFLATE} & The compression method used inside an \fname{.XLSX} archive.
\ccprod{} can decompress it but cannot produce it, which is why exported
\fname{.XLSX} files are large. \\

\textbf{Edit mode} & The state in which typing goes into a cell. The status
line says \msg{EDIT}. \\

\textbf{Error value} & A value such as \errv{\#DIV/0!} that a formula produces
when it cannot produce a number. Chapter~\ref{ch:messages} lists them. \\

\textbf{Formula} & A cell entry beginning with \lit{=}, evaluated to produce
the cell's value. \\

\textbf{Freeze} & Pinning the rows above the cursor and the columns to its left
so that they stay on screen while the rest scrolls.
Chapter~\ref{ch:moving}. \\

\textbf{Formula bar} & The second line of the screen, showing the current
cell's address and its actual contents. \\

\textbf{General} & The default number format: up to nine significant digits,
no trailing zeros, no leading zero below 1. \\

\textbf{Handle} & The bank-and-offset pair by which everything in banked RAM is
addressed, in place of a pointer. \\

\textbf{HostFS} & The emulator's way of presenting a directory on the host
computer as though it were a disk on unit 8. \\

\textbf{ISO mode} & The character mode \ccprod{} puts the machine into at
startup, which makes the keyboard produce ordinary ASCII. \\

\textbf{KERNAL} & The Commander X16's built-in operating system. \\

\textbf{MBF} & Microsoft Binary Format: the five-byte floating-point
representation \ccprod{} stores numbers in, and the one the machine's ROM
arithmetic uses. Appendix~\ref{app:numbers}. \\

\textbf{Number format} & How a cell's value is displayed --- General, Integer,
Decimal, Currency, Percent, Date, Time or Text. Chapter~\ref{ch:formats}. \\


\textbf{PETSCII} & The Commodore character encoding. \ccprod{} works in ASCII
internally and converts only where CMDR-DOS requires it, in file names. \\

\textbf{Range} & A rectangle of cells, written \cellref{A1:B9}. Only \fn{SUM},
\fn{AVERAGE}, \fn{MIN}, \fn{MAX} and \fn{COUNT} accept one. \\

\textbf{Ready mode} & The state in which keys move the cursor and run commands.
The status line says \msg{READY} or \msg{MODIF}. \\

\textbf{Recalculation} & Working out every formula's value again. It happens
after every change, over every worksheet, unless automatic recalculation has
been switched off. Chapter~\ref{ch:recalc}. \\

\textbf{Reference} & The name of a cell or a range inside a formula. \\

\textbf{Row} & A horizontal line of cells, numbered 1 to 65,535. \\

\textbf{Serial} & The number that stands for a date: 1 is 1~January~1900. \\

\textbf{Status line} & The second line from the bottom, showing the mode, the
current cell and the free memory. \\

\textbf{Style} & A stored combination of number format, decimal places, bold
and alignment. A cell names one by number. There are at most 128 in a
workbook. \\

\textbf{Tab} & The name of a worksheet, as shown on the tab line near the
bottom of the screen. \\

\textbf{Text pool} & The store holding every piece of text in the workbook.
Limited to 8,192 entries, and it does not share identical text. \\

\textbf{Used range} & The rectangle from \cellref{A1} to the last row and
column of a worksheet with anything in them. What \menupath{Export} writes and
what \menupath{Sort} covers. \\

\textbf{VERA} & The Commander X16's video chip. \ccprod{} writes to it
directly rather than going through the KERNAL, which is why the display is
fast and why \lit{-echo} cannot capture it. \\

\textbf{Workbook} & Everything in memory: up to sixteen worksheets, saved as
one \fname{.X16S} file. \\

\textbf{Worksheet} & One grid of cells, with its own name and its own tab. \\

\textbf{X16S} & \ccprod's own file format. Appendix~\ref{app:fileformat}. \\

\textbf{XLSX} & The Excel workbook format: a ZIP archive of XML documents.
\ccprod{} reads and writes it. Chapter~\ref{ch:importexport}. \\

\end{cctable}
